\documentclass[11pt,a4paper]{article}

% === 中文支持 ===
\usepackage{ctex}
\usepackage[T1]{fontenc}

% === 页面布局 ===
\usepackage[top=2.5cm, bottom=2.5cm, left=2.5cm, right=2.5cm]{geometry}

% === 表格 ===
\usepackage{booktabs}
\usepackage{array}
\usepackage{enumitem}
\usepackage{tabularx}

% === 数学符号 ===
\usepackage{amssymb}

% === 颜色 ===
\usepackage{xcolor}
\definecolor{primary}{HTML}{1a1a2e}
\definecolor{accent}{HTML}{3b82f6}
\definecolor{success}{HTML}{22c55e}
\definecolor{danger}{HTML}{ef4444}
\definecolor{warning}{HTML}{f59e0b}
\definecolor{graytext}{HTML}{6b7280}
\definecolor{progress}{HTML}{8b5cf6}

% === 超链接 ===
\usepackage[hidelinks]{hyperref}
\hypersetup{colorlinks=true, linkcolor=primary, urlcolor=accent}

% === 标题样式 ===
\usepackage{titlesec}
\titleformat{\section}{\Large\bfseries\color{primary}}{}{0em}{}[\titlerule]
\titleformat{\subsection}{\large\bfseries\color{primary}}{}{0em}{}
\titleformat{\subsubsection}{\normalsize\bfseries\color{primary}}{}{0em}{}

% === 页眉页脚 ===
\usepackage{fancyhdr}
\pagestyle{fancy}
\fancyhf{}
\fancyhead[L]{\small\color{graytext} 企业级 AI Agent 应用商店 -- 增量需求 DEMAND-016}
\fancyhead[R]{\small\color{graytext} v1.0}
\fancyfoot[C]{\small\color{graytext} \thepage}
\renewcommand{\headrulewidth}{0.4pt}

% === TikZ ===
\usepackage{tikz}
\usetikzlibrary{shapes.geometric, arrows.meta, positioning, calc, fit}
\usepackage{float}

\tikzset{
  arr/.style={thick, -Stealth},
  flowstep/.style={rectangle, rounded corners, draw=black, fill=white, thick,
    minimum width=2.4cm, minimum height=0.8cm, align=center, font=\small},
  rolehuman/.style={ellipse, draw=blue!60, fill=blue!12, thick,
    minimum width=2cm, minimum height=0.8cm, align=center, font=\small},
  rolesys/.style={rectangle, rounded corners, draw=gray!60, fill=gray!12, thick,
    minimum width=2cm, minimum height=0.8cm, align=center, font=\small},
  panel/.style={rectangle, rounded corners=4pt, draw=gray!50, fill=gray!4, thick,
    minimum width=3.5cm, minimum height=1.2cm, align=left, text width=3.2cm, font=\small},
}

% === 文档信息 ===
\title{
\vspace{-1em}
{\Huge\bfseries\color{primary} DEMAND-016}\\[0.6em]
{\Large\bfseries\color{accent} 用户设置与个人信息管理}\\[0.4em]
{\large 增量需求文档}
}
\author{
\textbf{PM AI}：万能产品小助手（OpenClaw）\\[0.3em]
\textbf{审阅}：魏子政
}
\date{2026-05-12 \quad 版本 v1.0}

\begin{document}

\maketitle
\thispagestyle{empty}

\begin{center}
\begin{tabular}{>{\bfseries}p{3cm} p{10cm}}
\toprule
需求编号 & DEMAND-016 \\
需求名称 & 用户设置与个人信息管理 \\
提出者 & 魏子政 \\
优先级 & \textcolor{warning}{P1} \\
状态 & 待开发 \\
关联需求 & D-015（邮箱认证）、D-008（用户模型） \\
技术栈 & Vue 3 + Element Plus + FastAPI + SQLAlchemy \\
\bottomrule
\end{tabular}
\end{center}

\vspace{1em}

\begin{abstract}
当前系统右上角用户下拉菜单仅有「退出登录」一项。用户无法查看或修改个人信息（用户名、邮箱、密码等），也无法进行偏好设置。本需求补充完整的用户自管理功能，包括个人资料编辑、密码修改、界面偏好设置。
\end{abstract}

\newpage
\tableofcontents
\newpage

% ============================================================
\section{背景与问题}
% ============================================================

\subsection{当前状态}

\begin{itemize}[nosep]
  \item 用户模型（\texttt{User}）字段：\texttt{username}, \texttt{hashed\_password}, \texttt{email}, \texttt{email\_verified}, \texttt{role}, \texttt{is\_active}
  \item 前端右上角下拉菜单：仅有「退出登录」
  \item 后端 API：仅有注册、登录、发送验证码、验证邮箱（\texttt{/api/auth/*}）
  \item 无「查看个人信息」接口，无「修改个人信息」接口，无「修改密码」接口
\end{itemize}

\subsection{用户痛点}

\begin{itemize}[nosep]
  \item 无法查看当前账号的注册信息（用户名、邮箱、角色）
  \item 无法修改密码
  \item 无法更新邮箱或重新发送验证邮件
  \item 无法切换界面语言偏好（虽然当前有 i18n，但无持久化设置）
\end{itemize}

% ============================================================
\section{需求描述}
% ============================================================

\subsection{总体目标}

在右上角下拉菜单中新增「个人设置」入口，点击后进入独立的用户设置页面，支持以下功能模块：

\begin{enumerate}[nosep]
  \item \textbf{个人资料}：查看和编辑用户名、邮箱
  \item \textbf{账号安全}：修改密码
  \item \textbf{界面偏好}：语言切换（中文/英文），主题切换（亮/暗）
\end{enumerate}

\subsection{页面布局设计}

设置页面采用经典的左侧 Tab + 右侧内容区的布局结构。

\begin{figure}[H]
\centering
\begin{tikzpicture}[transform shape, scale=0.85,
    node distance=0.4cm and 0.5cm,
  ]
  % 外框
  \draw[thick, rounded corners=6pt, fill=gray!3] (0,0) rectangle (12.5, 7.2);

  % 面包屑
  \node[anchor=north west, font=\small\color{graytext}] at (0.3, 6.9) {首页 > 个人设置};

  % 标题
  \node[anchor=north west, font=\Large\bfseries\color{primary}] at (0.3, 6.3) {个人设置};

  % 左侧 Tab 栏
  \draw[thick, rounded corners=3pt, fill=white] (0.3, 0.5) rectangle (3.3, 5.5);
  \node[anchor=north west, font=\small\bfseries] at (0.6, 5.3) {设置菜单};
  \draw[thick, fill=accent!10, rounded corners=2pt] (0.5, 4.2) rectangle (3.1, 4.7);
  \node[font=\small\bfseries\color{accent}] at (1.8, 4.45) {个人资料};
  \node[font=\small] at (1.8, 3.6) {账号安全};
  \node[font=\small] at (1.8, 2.8) {界面偏好};
  \node[font=\small] at (1.8, 2.0) {关于};

  % 右侧内容区 - 个人资料
  \draw[thick, rounded corners=3pt, fill=white] (3.8, 0.5) rectangle (12.2, 5.5);

  \node[anchor=north west, font=\large\bfseries\color{primary}] at (4.1, 5.3) {个人资料};

  % 头像区
  \draw[thick, circle, fill=accent!15, minimum size=1.2cm] (5.2, 4.3) circle (0.55);
  \node[font=\Large\bfseries\color{accent}] at (5.2, 4.3) {U};
  \node[font=\scriptsize\color{accent}, anchor=west] at (6.0, 4.3) {更换头像（V2）};

  % 表单字段
  \node[anchor=west, font=\small\color{graytext}] at (4.2, 3.5) {用户名};
  \draw[thick, rounded corners=2pt] (4.2, 2.9) rectangle (11.5, 3.3);
  \node[anchor=west, font=\small] at (4.4, 3.1) {current\_username};

  \node[anchor=west, font=\small\color{graytext}] at (4.2, 2.5) {邮箱地址};
  \draw[thick, rounded corners=2pt] (4.2, 1.9) rectangle (11.5, 2.3);
  \node[anchor=west, font=\small] at (4.4, 2.1) {user@example.com};
  \node[anchor=east, font=\scriptsize\color{success}] at (11.3, 2.1) {已验证};
  % 或 \node[...] {未验证，重新发送验证邮件}

  \node[anchor=west, font=\small\color{graytext}] at (4.2, 1.5) {账号角色};
  \draw[thick, rounded corners=2pt, fill=gray!8] (4.2, 0.9) rectangle (11.5, 1.3);
  \node[anchor=west, font=\small\color{graytext}] at (4.4, 1.1) {普通用户 / 管理员（只读）};

  % 保存按钮
  \draw[thick, rounded corners=3pt, fill=accent, text=white] (4.2, 0.6) rectangle (6.2, 0.85);
  \node[font=\small\bfseries\color{white}] at (5.2, 0.72) {保存修改};

\end{tikzpicture}
\caption{用户设置页面布局 — 个人资料 Tab}
\label{fig:settings-layout}
\end{figure}

\subsection{三个 Tab 的详细内容}

\subsubsection{Tab 1：个人资料}

\begin{table}[h]
\centering
\small
\begin{tabular}{llp{5cm}c}
\toprule
\textbf{字段} & \textbf{类型} & \textbf{说明} & \textbf{可编辑} \\
\midrule
用户头像 & 显示 & 首字母头像（V2 支持上传） & V2 \\
用户名 & 输入框 & 2-64 字符，唯一性校验 & 是 \\
邮箱地址 & 输入框 & 标准邮箱格式，唯一性校验 & 是 \\
邮箱验证状态 & 标签 & 已验证（绿）/ 未验证（橙），未验证时显示「重新发送验证邮件」链接 & 只读 \\
账号角色 & 标签 & 管理员 / 普通用户 & 只读 \\
注册时间 & 文本 & 格式化日期 & 只读 \\
\bottomrule
\end{tabular}
\end{table}

\textbf{交互规则}：
\begin{itemize}[nosep]
  \item 修改用户名：实时校验唯一性（防抖 500ms），保存时二次校验
  \item 修改邮箱：保存后触发重新验证流程（发送验证码邮件），邮箱验证状态重置为「未验证」
  \item 保存成功后刷新页面状态，右上角头像文字同步更新
\end{itemize}

\subsubsection{Tab 2：账号安全}

\begin{table}[h]
\centering
\small
\begin{tabular}{llp{6cm}}
\toprule
\textbf{字段} & \textbf{类型} & \textbf{说明} \\
\midrule
当前密码 & 密码框 & 必填，验证身份 \\
新密码 & 密码框 & 8-64 字符，需包含字母+数字 \\
确认新密码 & 密码框 & 必须与新密码一致 \\
\bottomrule
\end{tabular}
\end{table}

\textbf{交互规则}：
\begin{itemize}[nosep]
  \item 修改密码前必须验证当前密码
  \item 新密码强度校验：至少 8 位，包含字母和数字
  \item 修改成功后自动刷新 Token（相当于重新登录），不跳转到登录页
  \item 连续输错 5 次当前密码，锁定 15 分钟（V2）
\end{itemize}

\subsubsection{Tab 3：界面偏好}

\begin{table}[h]
\centering
\small
\begin{tabular}{llp{5cm}}
\toprule
\textbf{设置项} & \textbf{类型} & \textbf{说明} \\
\midrule
界面语言 & 单选 & 中文 / English，即时生效，持久化到 localStorage \\
主题模式 & 开关 & 亮色 / 暗色（当前已有暗色主题变量，需持久化开关状态） \\
\bottomrule
\end{tabular}
\end{table}

\textbf{交互规则}：
\begin{itemize}[nosep]
  \item 语言切换即时生效，刷新页面后保持
  \item 主题切换即时生效，刷新页面后保持
  \item 偏好设置保存在 localStorage（MVP 阶段），V2 可同步到后端用户配置
\end{itemize}

% ============================================================
\section{后端 API 设计}
% ============================================================

\subsection{新增接口}

\begin{table}[h]
\centering
\small
\begin{tabular}{llp{4cm}p{4cm}}
\toprule
\textbf{方法} & \textbf{路径} & \textbf{说明} & \textbf{请求/响应} \\
\midrule
GET & \texttt{/api/auth/me} & 获取当前用户信息 & 响应：\texttt{\{username, email, email\_verified, role, created\_at\}} \\
PUT & \texttt{/api/auth/me} & 修改个人资料 & 请求：\texttt{\{username?, email?\}}，响应：更新后的用户信息 \\
PUT & \texttt{/api/auth/change-password} & 修改密码 & 请求：\texttt{\{current\_password, new\_password\}}，响应：新 Token \\
\bottomrule
\end{tabular}
\end{table}

\textbf{接口细节}：

\texttt{GET /api/auth/me}：
\begin{itemize}[nosep]
  \item 需要 Bearer Token
  \item 返回 \texttt{UserPublic}（复用注册时的响应模型）+ \texttt{created\_at}
  \item 不返回 \texttt{hashed\_password}、\texttt{is\_active} 等敏感字段
\end{itemize}

\texttt{PUT /api/auth/me}：
\begin{itemize}[nosep]
  \item 字段均为可选（partial update），只传需要修改的字段
  \item 修改用户名：校验唯一性（排除自身），2-64 字符
  \item 修改邮箱：校验唯一性（排除自身）、格式校验，修改后 \texttt{email\_verified} 重置为 false
  \item 修改邮箱后自动触发验证码发送（如 SMTP 已配置）
\end{itemize}

\texttt{PUT /api/auth/change-password}：
\begin{itemize}[nosep]
  \item 验证 \texttt{current\_password} 正确性
  \item 新密码强度校验
  \item 返回新 Token（旧 Token 失效）
\end{itemize}

\subsection{前端路由}

新增路由：
\begin{itemize}[nosep]
  \item \texttt{/settings} $\to$ \texttt{SettingsView.vue}（需登录）
\end{itemize}

右上角下拉菜单更新：
\begin{itemize}[nosep]
  \item 新增「个人设置」菜单项（排在「退出登录」之前）
  \item 点击跳转 \texttt{/settings}
\end{itemize}

% ============================================================
\section{任务分解}
% ============================================================

\begin{table}[h]
\centering
\small
\begin{tabular}{clp{5cm}c}
\toprule
\textbf{序号} & \textbf{任务} & \textbf{说明} & \textbf{预估} \\
\midrule
1 & 后端：用户信息查询接口 & \texttt{GET /api/auth/me}，返回 UserPublic + created\_at & 0.5h \\
2 & 后端：用户信息修改接口 & \texttt{PUT /api/auth/me}，支持用户名/邮箱修改 & 0.5h \\
3 & 后端：修改密码接口 & \texttt{PUT /api/auth/change-password} & 0.5h \\
4 & 前端：SettingsView 页面 & 三 Tab 布局，表单交互 & 2h \\
5 & 前端：API 函数 \& i18n & 前端 API 封装 + 中英文翻译 & 0.5h \\
6 & 前端：路由 \& 导航菜单 & 新增 /settings 路由，下拉菜单添加入口 & 0.5h \\
7 & 联调 \& 测试 & 接口联调、表单校验、边界测试 & 1h \\
\midrule
 & & \textbf{合计} & \textbf{5.5h} \\
\bottomrule
\end{tabular}
\end{table}

% ============================================================
\section{验收标准}
% ============================================================

\begin{enumerate}[nosep]
  \item[\textcolor{success}{\checkmark}] 右上角下拉菜单有「个人设置」和「退出登录」两个选项
  \item[\textcolor{graytext}{$\square$}] 点击「个人设置」进入设置页面，左 Tab 右内容布局
  \item[\textcolor{graytext}{$\square$}] 「个人资料」Tab 可查看用户名、邮箱、角色、注册时间
  \item[\textcolor{graytext}{$\square$}] 修改用户名：保存成功，唯一性冲突时提示
  \item[\textcolor{graytext}{$\square$}] 修改邮箱：保存成功，验证状态重置为「未验证」
  \item[\textcolor{graytext}{$\square$}] 「账号安全」Tab：输入当前密码 + 新密码，修改成功后自动刷新 Token
  \item[\textcolor{graytext}{$\square$}] 「界面偏好」Tab：语言/主题切换即时生效，刷新后保持
  \item[\textcolor{graytext}{$\square$}] 全部表单有中文/英文 i18n
  \item[\textcolor{graytext}{$\square$}] 未登录状态访问 \texttt{/settings} 跳转登录页
\end{enumerate}

% ============================================================
\section{V2 远期规划}
% ============================================================

\begin{itemize}[nosep]
  \item 用户头像上传（MinIO 存储）
  \item 修改密码连续错误锁定机制
  \item 偏好设置同步到后端（跨设备一致）
  \item 通知设置（邮件通知开关）
  \item API Token 管理（生成/吊销个人访问令牌）
\end{itemize}

% ============================================================
\section*{更新记录}
% ============================================================

\begin{tabular}{lll}
\toprule
版本 & 日期 & 变更内容 \\
\midrule
v1.0 & 2026-05-12 & 初始版本 \\
\bottomrule
\end{tabular}

\end{document}
